\section{An exact global unit criterion}
\label{sec:cm-global-unit-criterion}

Fix a prime \(p\geq5\), write \(h=(p-1)/2\), choose a primitive
residue \(g\) modulo \(p\), and put
\[
 F=\mathbb Q(\zeta_p),\quad F^+=\mathbb Q(\zeta_p+\zeta_p^{-1}),
 \quad \mathcal O=\mathbb Z[\zeta_p]\otimes\mathbb Z_2,\quad
 s=[g],\quad c=s^h.
\]
The algebra \(\mathcal O\) may be a product. All its factors are
included in the following groups:
\[
 T=\{t\in1+2\mathcal O:t\,c(t)=1\},\qquad P=(s-1)T.
\]
We use additive notation for operators on these multiplicative
abelian groups. Let \(Q=Q_p(s)\) be the integral reflex operator and
assume \(Q(1)\) is odd. Set
\[
 \mathscr M=(\mathcal O/2\mathcal O)^c,\qquad
 \beta_k=\overline\zeta_p^{\,g^k}+
                \overline\zeta_p^{-g^k}\quad(0\leq k<h).
\]
These elements form an \(\mathbb F_2\)-basis on which \(s\) acts cyclically. The hypothesis
needed here is the natural finite isomorphism
\begin{equation}
 \theta:T/QT\xrightarrow{\ \sim\ }
               \mathscr M/\overline Q\mathscr M,\qquad
 \theta(t)=(t-1)/2\bmod2\mathcal O.
 \label{eq:uniform-natural-layer}
\end{equation}
An equality of rational ranks is insufficient for this hypothesis.

\begin{lemma}
\label{lem:uniform-operator-inverses}
Under \eqref{eq:uniform-natural-layer}, inclusion identifies
\(P/QP\) with \(T/QT\). On this quotient both \(s-1\) and \(1+[2]\)
are automorphisms.
\end{lemma}
\begin{proof}
The operator \(Q\) acts on \(T/P\) by the odd scalar \(Q(1)\).
It is therefore an automorphism of this pro-\(2\) group. Hence
\(T=P+QT\), and \(Qt\in P\) implies \(t\in P\); consequently
\(QT\cap P=QP\).

The map \(\theta:T\to\mathscr M\) is surjective. Indeed, on each
unramified factor complex conjugation is either the nontrivial
automorphism of an unramified quadratic extension or exchanges
two factors. Thus \(1+c:\mathcal O/2\mathcal O\to\mathscr M\) is
surjective. If \(w+c(w)\equiv v\pmod2\), then
\((1+2w)/(1+2c(w))\) belongs to \(T\) and has image \(v\).
Also \(T[2]\) consists exactly of componentwise signs, equal on
the factors exchanged by \(c\). Since \([2]\) on
\(\mathcal O/2\mathcal O\) is absolute Frobenius, this gives
\begin{equation}
 T[2]\xrightarrow{\ \theta\ }\ker(1+[2]:\mathscr M\to\mathscr M)
 \quad\text{isomorphically}.
 \label{eq:uniform-sign-kernel}
\end{equation}
Let \(b=\dim_{\mathbb F_2}(\mathscr M/\overline Q\mathscr M)\).
The integral anti-invariant lattice \(\mathcal O^-\) has basis
\(\zeta_p^{g^k}-\zeta_p^{-g^k}\), and is
\(\mathbb Z_2[s]/(s^h+1)\). Its reduction is \(\mathscr M\).
The logarithm identifies \(T\otimes\mathbb Q_2\) with
\(\mathcal O^-\otimes\mathbb Q_2\). The Smith normal form therefore
gives
\[
 e:=v_2\det(Q\mid T\otimes\mathbb Q_2)\geq b.
\]
The exact index on the finitely generated abelian pro-\(2\) group
\(T\) is
\begin{equation}
 [T:QT]=|\ker(Q:T[2]\to T[2])|\,2^e.
 \label{eq:uniform-sign-index}
\end{equation}
To see this formula, pass to the torsion-free quotient of \(T\).
The induced endomorphism is injective, its cokernel has order
\(2^e\), and the remaining cokernel is the cokernel of \(Q\)
on the finite torsion group. Kernel and cokernel there have equal
order. Hypothesis \eqref{eq:uniform-natural-layer} forces \(e=b\)
and forces \(Q\) to be invertible on \(T[2]\).

Put \(D=1+[2]\). In the Fitting decomposition
\(\mathscr M=\mathscr M_0\oplus\mathscr M_1\) for \(D\),
the operator \(D\) is nilpotent on the first summand and invertible
on the second. The operator \(\overline Q\) commutes with \(D\).
A nonzero kernel vector of \(\overline Q\) in \(\mathscr M_0\)
would have a last nonzero \(D\)-iterate in
\(\ker D\cap\ker\overline Q\), contrary to
\eqref{eq:uniform-sign-kernel}. Thus the first summand contributes
nothing to the cokernel of \(\overline Q\), and \(D\) is invertible
on that cokernel. Finally, \(\overline Q(1)=1\), so a Bezout
identity for \(X-1\) and \(\overline Q(X)\) makes \(s-1\)
invertible there as well. No semisimplicity is used.
\end{proof}

Write
\[
 \mathscr C=P/QP=T/QT=\mathscr M/\overline Q\mathscr M,
 \qquad \pi:\mathscr M\to\mathscr C.
\]
Let \(\mathscr E=\mathcal O_{F^+}^{\times}\) be the full real
unit group. Define the linear square-class map
\begin{equation}
 \mathscr K:\mathscr C^\vee\longrightarrow
       \mathscr E/(\{\pm1\}\mathscr E^2),\qquad
 \ell\longmapsto[u_\ell],\qquad
 u_\ell=\prod_{k=0}^{h-1}
  (\zeta_p^{g^k}+\zeta_p^{-g^k})^{\ell(\pi\beta_k)}.
 \label{eq:uniform-real-unit-map}
\end{equation}
Each factor is a unit: up to a root of unity it is
\(1+\zeta_p^{2g^k}\), whose norm is \(\Phi_p(-1)=1\).
Square roots in \eqref{eq:uniform-real-unit-map} may be arbitrary
units of \(F^+\).

\begin{theorem}
\label{thm:uniform-arithmetic-cm-image}
Let \(R_\Lambda\) and \(V_\Lambda\) be obtained from
\(R_{\rm quad}\) and \(V_{\rm quad}\) by imposing any collection
of the scalar kernels \(h_q\) with \(q\ne p\).
Under \eqref{eq:uniform-natural-layer}, their arithmetic images are
\begin{equation}
 h_p(R_\Lambda\cap H)=h_p(V_\Lambda\cap H)
  =\{v\in P:\ell(v\bmod QP)=0\text{ for every }
                         \ell\in\ker\mathscr K\}.
 \label{eq:uniform-arithmetic-image}
\end{equation}
In particular,
\[
 [P:h_p(R_\Lambda\cap H)]=2^{\dim\ker\mathscr K},\qquad
 \dim_{\mathbb F_2}\bigl(h_p(R_\Lambda\cap H)/QP\bigr)
                      =\dim\operatorname{im}\mathscr K.
\]
Every value in \eqref{eq:uniform-arithmetic-image} has a witness
in \(V_{\rm quad}\cap H\) with all other prime scalar coordinates
equal to \(1\).
\end{theorem}

\begin{proof}[Proof of Theorem~\ref{thm:uniform-arithmetic-cm-image}]
First, a real unit \(u\in\mathscr E\) is a square in \(F(i)\)
if and only if \(u\in\{\pm1\}\mathscr E^2\).
Indeed \(F(i)/F^+\) is biquadratic: \(F\) is unramified at \(2\),
so it does not contain \(i\). Of its three quadratic subfields,
\(F\) and the third, real subfield are ramified at the prime
above \(p\); \(F^+(i)\) is unramified there. Adjoining a square
root of a unit is unramified at every odd prime. Thus a nontrivial
\(F^+(\sqrt u)\) inside \(F(i)\) must be \(F^+(i)\).
Writing a square root as \(a+bi\) gives \(ab=0\), so \(u=a^2\)
or \(u=-b^2\). The square root is a unit, by valuations.

Let \(\rho:G_F\to\mathcal O^\times\) be the scalar of the compact-homology
Tate representation, and remove the odd Teichmuller factor in
every component to obtain its principal projection
\(\operatorname{pr}_U\rho\). For \(b\in G_{F(i)}\), set
\[
 v_b=\bigl(\operatorname{pr}_U\rho(b)-1\bigr)/2\bmod2\mathcal O.
\]
Polarization and \(\lambda_2(b)\equiv1\pmod4\) imply
\(v_b\in\mathscr M\). This is a continuous additive character.
For \(\ell\in\mathscr C^\vee\), define
\[
 \mathscr D(\ell)=
 \left[\prod_{k=0}^{h-1}
   (1-\zeta_p^{g^k})^{\ell(\pi\beta_k)}\right]
          \in F(i)^\times/(F(i)^\times)^2.
\]
We claim
\begin{equation}
 \ell(\pi v_b)=\kappa_{\mathscr D(\ell)}(b),
 \label{eq:uniform-jacobi-kummer}
\end{equation}
where \(\kappa\) is the quadratic Kummer character.

It suffices to check Frobenius elements at primes of \(F(i)\)
away from \(2p\). Let \(q\) be the full residue-field cardinality,
\(m=(q-1)/p\), and let \(\chi\) be the residue character with
\(\chi(x)\equiv x^m\) at that prime and \(\chi(0)=0\).
Then \(q\equiv1\pmod{4p}\) and \(m\) is even. The
character comparison in Theorem~8.1, Remark~8.3, and
Corollary~8.4 of \cite{Otsubo2016} applies over the number field
\(F(i)\), with auxiliary residue degree \(1\), and gives the
negative Jacobi sum
\[
 j=-\sum_{x\in\mathbb F_q}\chi(x)\chi(1-x).
\]
The substitution \(y=2x-1\) gives
\[
 \chi(4)j=-\sum_{y\in\mathbb F_q}\chi(1-y^2)
        =-1-2\sum_{y\in\mathbb F_q^\times/\{\pm1\}}
                                      \chi(1-y^2).
\]
This element is \(1\) modulo \(2\), so \(\chi(4)^{-1}\) is
exactly the odd Teichmuller factor of \(j\). Therefore
\[
 v_{\operatorname{Frob}}\equiv
 1+\sum_{y\in\mathbb F_q^\times/\{\pm1\}}\chi(1-y^2)
                                                   \pmod2.
\]
For \(a\ne0\), let \(n_a\) count the pairs contributing
\(\zeta_p^a\). Equivalently, it counts roots of \(z^m=\zeta_p^a\)
for which \(1-z\) is a nonzero square. All \(m\) roots lie in
\(\mathbb F_q\), none is \(1\), and
\[
 \prod_{z^m=\zeta_p^a}(1-z)=1-\zeta_p^a.
\]
The quadratic residue sign of this product is
\((-1)^{m-n_a}=(-1)^{n_a}\). Since \(-1\) and \(\zeta_p\)
are squares in \(\mathbb F_q\), opposite exponents have the
same parity. Thus
\[
 v_{\operatorname{Frob}}=
 \epsilon\cdot1+\sum_{k=0}^{h-1}
     \kappa_{1-\zeta_p^{g^k}}(\operatorname{Frob})\beta_k
       \quad\text{in }\mathscr M.
\]
The constant term is killed by \(\pi\), because
\(\overline Q(1)=1\). Applying \(\ell\pi\) proves
\eqref{eq:uniform-jacobi-kummer} on Frobenius elements;
Chebotarev proves it on all of \(G_{F(i)}\).

Put \(d=[2]\) as an operator, and \(L_0=d^{-2}(1+d)\).
The identity
\[
 \zeta_p^a+\zeta_p^{-a}
  =\zeta_p^{-a}\frac{1-\zeta_p^{4a}}{1-\zeta_p^{2a}}
\]
gives, in square classes of \(F(i)\),
\begin{equation}
 [u_\ell]=\mathscr D(\ell\circ L_0).
 \label{eq:uniform-dual-reindexing}
\end{equation}
The inverse exponents in this formula are essential:
the coefficient at \(1-\zeta_p^a\) is
\(\ell(\pi\beta_{a/4})+\ell(\pi\beta_{a/2})\).
Replacing an exponent by its negative changes the factor
only by a square in \(F(i)\).

The kernel of \(\mathscr D\) is stable under precomposition
by \(s\), since
\(\mathscr D(\ell\circ s)=s^{-1}\mathscr D(\ell)\).
It is therefore stable under every polynomial in \(s\).
Lemma~\ref{lem:uniform-operator-inverses} makes \(L_0\)
and \(s-1\) invertible on \(\mathscr C\). Their dual actions
preserve this finite kernel and are bijective on it.
The real-unit square test and \eqref{eq:uniform-dual-reindexing}
consequently give
\[
 \ker\mathscr K=\ker\mathscr D.
\]
Finite-dimensional duality in \eqref{eq:uniform-jacobi-kummer}
now shows that
\begin{equation}
 \{\pi v_b:b\in G_{F(i)}\}
       =W:=(\ker\mathscr K)^\perp\subseteq\mathscr C.
 \label{eq:uniform-image-annihilator}
\end{equation}
The subspace \(W\) is stable under \(s\) and under
\((s-1)^{-1}\).

Every element of \(R_\Lambda\cap H\) or \(V_\Lambda\cap H\)
fixes \(F(i)\), and its principal Tate scalar is \(h_p\).
This proves the upper bound in
\eqref{eq:uniform-arithmetic-image}. For the reverse bound,
use the established mixed-inertia construction. Choose primitive
\(p\)-inertia \(a\), and take \(b\in G_{F(i)}\). The word
\[
 w_b=[a,[a,b]^{e_2}]
\]
belongs to \(V_{\rm quad}\cap H\), and its scalar coordinates
in every degree \(q\ne p\) are \(1\). In degree \(p\),
the scalar of \([a,b]\) is \((s-1)\rho(b)\).
The pro-\(2\) power removes its odd part; the outer commutator
applies \(s-1\) once more. Hence
\[
 h_p(w_b)\bmod QP=(s-1)^2\pi v_b.
\]
By \eqref{eq:uniform-image-annihilator}, these words realize
all of \(W\) modulo \(QP\). The local norm construction realizes
every element of \(QP\), again with every other coordinate
equal to \(1\). Multiplication gives every value in the claimed
preimage, with all other coordinates equal to \(1\).
Such witnesses survive any collection, finite or infinite,
of other-degree scalar kernels. This proves the exact image
and both dimension formulas.
\end{proof}

